% ============================================================
%  Section 5 -- Design patterns
% ============================================================

\chapter{Patrons de conception et modèles d'analyse}

\section{Cas d'utilisation}

Le système vise plusieurs acteurs : l'étudiant ou développeur qui lance les
benchmarks, l'encadrant qui analyse les résultats, le support IT cible, les
agents distants A2A et le pipeline CI. Les principaux cas d'utilisation sont
representes ci-dessous.

\plantumldiagram{use_cases}{Diagramme des cas d'utilisation de \bibopsname{}.}{fig:use-cases}

\section{Diagramme de classes d'analyse}

Le diagramme d'analyse se concentre sur les concepts metier : ticket, agent,
outil, trace, évaluateur, résultat de benchmark et politique composite. Il ne
cherche pas à représenter chaque classe Python, mais les responsabilités qui
structurent la solution.

\plantumldiagram{classes_analyse}{Diagramme de classes d'analyse.}{fig:classes-analyse}

\section{Patrons identifiés}

Plusieurs patrons de conception apparaissent dans le dépôt.

\begin{table}[H]
\centering
\begin{tabularx}{\textwidth}{lY}
\toprule
\textbf{Patron} & \textbf{Usage dans le projet} \\
\midrule
Adapter & Unifie les appels vers l'agent IT, les agents A2A et les APIs OpenAI-compatible. \\
Strategy & Permet de changer d'évaluateur ou de politique de score sans réécrire les runners. \\
Registry & Centralise les évaluateurs disponibles et les rend extensibles. \\
Facade & La CLI \code{bibops} masque la complexite des modules internes. \\
Observer & La Racing Arena diffuse la télémétrie SSE à plusieurs équipes abonnées. \\
Template Method & Les runners de benchmark partagent la sequence charger, exécuter, évaluer, exporter. \\
\bottomrule
\end{tabularx}
\caption{Patrons de conception principaux.}
\label{tab:patterns}
\end{table}

\plantumldiagram{design_patterns}{Vue PlantUML des patrons de conception.}{fig:design-patterns}

\section{Choix de conception}

Le choix le plus important est la séparation entre le raisonnement du modèle et
les sources verifiees. Les outils retournent des données bornees, les décisions
de l'agent sont structurées et les traces conservent les appels. Cette
séparation rend l'évaluation possible : lorsqu'une réponse est mauvaise, on peut
identifier si l'erreur vient du routage, de l'outil, du modèle ou du juge.

Le second choix est le recours à des contrats simples. SQLite suffit pour un
statut serveur, JSON suffit pour une knowledge base de prototype, ChromaDB
permet le RAG sans infrastructure lourde, et Typer fournit une CLI claire. Cette
sobriété réduit le nombre de composants à maintenir et facilite les tests.
